\documentclass[11pt]{article}
\usepackage[margin=1.05in]{geometry}
\usepackage{amsmath,amssymb}
\usepackage{microtype}
\usepackage{xcolor}
\usepackage[colorlinks=true,linkcolor=blue!55!black,urlcolor=blue!55!black]{hyperref}
\usepackage{enumitem}
\setlist{topsep=2pt,itemsep=1pt}
\newcommand{\eps}{\varepsilon}
\newcommand{\dd}{\mathrm{d}}
\newcommand{\Q}{\mathbb{Q}}
\newcommand{\Z}{\mathbb{Z}}
\newcommand{\Fp}{\mathbb{F}_p}

\title{Differential Equations for Master Integrals\\[0.4ex]
\large A concise overview}
\date{August 24, 2026}

\begin{document}
\maketitle

\noindent
These notes summarize a longer treatment.  The setting is the
double-real channel of $pp \to h + X$ at NNLO, in $d = 4-2\eps$
dimensions.  Reverse unitarity replaces each on-shell delta function
by the discontinuity of a propagator, on positive-energy
configurations,
\begin{equation}
  \delta^{(+)}(p^2) \;=\; \frac{1}{2\pi i}
  \left[\frac{1}{p^2-i0} - \frac{1}{p^2+i0}\right] ,
\end{equation}
so that the phase-space integrals become cut two-loop integrals: they
obey the ordinary integration-by-parts identities, a cut line raised
to power $a_i$ standing for
$\big[(-1)^{a_i-1}/(a_i-1)!\big]\,\delta^{(a_i-1)}$ and an integral
with a cancelled cut line vanishing.  Reduction leaves $347$ master
integrals.  A master carries the overall scale by its mass dimension,
$I_m = (-s)^{Ld/2-a_m}\hat I_m$ with $L$ loops and propagator powers
summing to $a_m$, and the equations below govern the dimensionless
$\hat I_m$.  For the partonic process
$p_1 + p_2 \to p_3 + k_1 + k_2$ with all partons massless, and
$s = (p_1+p_2)^2$, $t = (p_1-p_3)^2$, $u = (p_2-p_3)^2$, these depend
on the two ratios $v = -t/s$ and $w = -u/s$ alone.  Since
$s+t+u = (p_1+p_2-p_3)^2 = M^2$ is the squared invariant mass of the
unobserved pair, $1-v-w = M^2/s \ge 0$, while $v, w > 0$: the
physical region is the triangle $0<v$, $0<w$, $v+w<1$.  In the
partonic centre-of-mass frame
$v = (E_3/\sqrt s)(1-\cos\theta)$ and
$w = (E_3/\sqrt s)(1+\cos\theta)$, so $v \to 0$ and $w \to 0$ are the
collinear limits and $v+w \to 1$ the Born configuration
$E_3 = \sqrt s/2$, $M^2 = 0$, at which the radiation is unresolved:
the soft edge.  The problem is the analytic evaluation of the
masters: each as a Laurent series in $\eps$ whose coefficients are
explicit transcendental functions with exact constants.

\section{Differential equations}
A master depends on the external momenta only through their
invariants $s_{kl}$, so the chain rule
$p_j^\mu\,\partial_{p_i^\mu} I = \sum_{k\le l}
\big(p_j^\mu\,\partial s_{kl}/\partial p_i^\mu\big)\,
\partial I/\partial s_{kl}$ is a linear system for the
$\partial I/\partial s_{kl}$, invertible wherever the Gram
determinant of the external momenta is non-zero; inverting it
expresses $\partial_v$ and $\partial_w$ through the operators
$p_j^\mu\,\partial/\partial p_i^\mu$.  Taken under the integral sign
these raise propagator powers and produce numerators which, the
family being complete, are linear combinations of its inverse
propagators and of invariants; a derivative of a master is therefore
again a combination of integrals of the same family, and
integration by parts brings it back to the masters of the family and
its subsectors.  Each family therefore satisfies a closed first-order
system in both variables at once,
\begin{equation}
  \dd \vec I = A(v,w;\eps)\,\vec I,
  \qquad A = A_v\,\dd v + A_w\,\dd w ,
\end{equation}
with rational matrices --- a \emph{Pfaffian system}, in the
terminology of the literature --- subject to the integrability
(flatness) condition $\dd A = A\wedge A$, verified exactly.  The
families themselves are already fixed by the reduction: propagator
sets are identified modulo \emph{affine} redefinitions of the loop
momenta, a linear map composed with a shift,
$k_\ell \mapsto \sum_m M_{\ell m} k_m + \sum_r c_{\ell r} p_r$ with
$M \in \mathrm{GL}(L,\Z)$, hence $\lvert\det M\rvert = 1$ and the
measure invariant, that map cut lines to cut lines.  The Symanzik
polynomials $(U,F)$ depend only on the graph and on the assignment of
the external momenta to its vertices, not on the routing of the loop
momenta, so equivalent families share them up to a permutation of the
Feynman parameters; a coincidence of their normal orderings (Pak's
algorithm) nominates a candidate, and the explicit momentum shift,
constructed from the linear conditions
$\widetilde D_i(Mk + c\cdot p) = D_{\sigma(i)}(k)$ and applied,
proves it.  Because differentiation never leaves a family, the $347$
masters split into $91$ separate Pfaffian systems, of dimensions
between $2$
and $47$ (the dimensions sum to more than $347$ because each system
carries the subsector masters it couples to, counted once in the
family that owns them).  The systems are separate but not
independent: a subsector master shared by two of them is one and the
same function, and they are solved in an order compatible with the
subsector relation.  Ordered by sectors ---
a sector is the set of propagators actually present --- each
connection is block lower triangular: diagonal blocks couple the
masters of one sector group, and the entries below the diagonal
feed subsector solutions upward.

\section{Equivalence of blocks}
Form the directed graph on the masters of a system with an edge
$i \to j$ whenever $I_i$ occurs in the derivative of $I_j$.  A
\emph{block} is a strongly connected component of that graph,
equivalently a diagonal block of the triangular structure.  The $91$
systems contain $1119$ blocks.  Two blocks with connections
$b^{(1)}$, $b^{(2)}$ are
\emph{equivalent} when there is a constant permutation matrix $P$
with
\begin{equation}
  b^{(2)}(v,w) = P^{-1}\, b^{(1)}(v,w)\, P
  \qquad\text{or}\qquad
  b^{(2)}(v,w) = P^{-1}\, b^{(1)}(w,v)\, P ,
\end{equation}
the second option being the exchange $p_1 \leftrightarrow p_2$ of the
incoming partons --- an exact matrix identity, never an analogy.
Since $P$ is constant it carries no derivative term, so the
equivalence is decided entry by entry.  Only $173$ equivalence
classes occur, because the same subtopologies recur inside many
parent families.
The analysis is therefore done once for each class; a family is
recovered afterwards from the classes it contains, and the
identification is checked again at that point.

\section{The canonical form}
A system is in \emph{canonical form} when
\begin{equation}
  \dd \vec F = \eps \sum_a R_a\, \dd\log \varphi_a\;\vec F ,
\end{equation}
with constant residue matrices $R_a$ and $\eps$-free \emph{letters}
$\varphi_a$ (the \emph{alphabet}).  For two blocks in this form the
equivalence of the preceding section reduces to a statement about the
residues alone, $R^{(2)}_a = P^{-1}R^{(1)}_{\sigma(a)}P$, with
$\sigma$ the relabeling of the letters induced by the exchange.  The
Laurent expansion decouples order by order,
$\dd\vec F^{(k)} = \big(\sum_a R_a\,\dd\log\varphi_a\big)\vec
F^{(k-1)}$: each order is a quadrature of the previous one, and the
solution is a path-ordered exponential whose order-$\eps^k$ term is
a $\Q$-linear combination of $k$-fold iterated integrals over the
alphabet, contracted with the boundary vector --- of uniform
transcendental weight $k$, and, on segments where every letter is
linear in the path parameter, generalized polylogarithms.  Take the
alphabet
$\{w, 1-w\}$ with
\begin{equation}
  R_w = \begin{pmatrix} 0 & 0\\ 0 & -1\end{pmatrix},
  \qquad
  R_{1-w} = \begin{pmatrix} 0 & 0\\ 1 & 0\end{pmatrix},
  \qquad
  \vec F \to (1,0)^{\mathsf T} \ \text{ at } \ w \to 0 .
\end{equation}
The first component stays constant, and the recursion gives
$F_2^{(1)} = \int_0^w \dd\log(1-w') = \log(1-w)$ and then
$F_2^{(2)} = -\int_0^w \dd\log w'\,\log(1-w')
= \operatorname{Li}_2(w)$: the dilogarithm as the twofold iterated
integral of the word $(w,\,1{-}w)$, and nothing more elaborate is
involved.  All singular behaviour is manifest: in the finite plane
the connection has simple poles exactly at the letters' zeros, and
the residue matrix $R$ there prescribes the local solutions
$H(w)\,w^{\eps R}\vec c$ with $H$ analytic and $H(0)=1$, that is
modes $w^{\,n+a\eps}(\log w)^j$ --- exponents from $\eps$ times the
eigenvalues of $R$, logarithms at fixed $\eps$ from its Jordan
structure --- known before any integration.  Two provisos: the normal
form requires that no two exponents differ by a nonzero integer,
automatic for generic $\eps$; and after expansion in $\eps$
logarithms appear at every order in any case, since
$w^{\eps R} = \sum_m (\eps\log w)^m R^m/m!$, whether or not $R$ is
diagonalizable.

\section{Construction of the transformation}
Under a change of basis $\vec I = T\vec F$ the connection transforms
as $A \mapsto T^{-1}AT - T^{-1}\dd T$, and the triangular structure
decides the order in which the blocks are treated.  A diagonal block
never feels what lies below it, so each block is brought to canonical
form on its own, once per class.  Either one posits a rational $T$ of
bounded degree, over the alphabet supplied by the irreducible factors
of the denominators of $A$ and of the discriminants in its residue
eigenvalues, and matches coefficients in
$\dd T = A\,T - \eps\,T\sum_a R_a\,\dd\log\varphi_a$ --- a system
bilinear in $T$ and the residues, so one fixes candidate residues and
solves linearly for $T$ in each case; or one reduces higher poles by
Moser's algorithm, which decides by a rank condition whether a
rational transformation lowers the pole, and normalizes the residue
eigenvalues by balances
\begin{equation}
  T = \bar P + \frac{v-v_1}{v-v_2}\,P ,
  \qquad \bar P = 1-P ,
  \qquad P = \frac{u\,\eta^{\mathsf T}}{\eta^{\mathsf T}u} ,
\end{equation}
with $\eta$ a left eigenvector of the residue at $v_1$, eigenvalue
$\lambda$, and $u$ a right eigenvector of the residue at $v_2$,
eigenvalue $\mu$, and $\eta^{\mathsf T}u \ne 0$.  Both conditions are
needed: $PA_{(1)}\bar P = 0$ keeps the transformed connection free of
a double pole at $v_1$ and $\bar P A_{(2)}P = 0$ does the same at
$v_2$, and a spectral projector of one residue satisfies only the
first.  The balance sends $\lambda\mapsto\lambda-1$ and
$\mu\mapsto\mu+1$ and leaves the exponents at every other point
unchanged, the residue at $v_k$ being conjugated by the constant
$T(v_k)$ and the one at infinity untouched since $T(\infty)=1$; it
must act at two points, since the sum of all residues including the
one at infinity, $-\sum_k A_k$, vanishes identically, so no
transformation changes an exponent at a single point.  Repeated
balances remove the integer parts $n$ from all eigenvalues
$n + a\eps$ --- each admissible step decreases the non-negative
integer $\sum|n|$ by two, so the sequence cannot continue
indefinitely; what can fail is the existence of an admissible pair with
$\eta^{\mathsf T}u \neq 0$ at some step.  A constant transformation
then factors $\eps$ out of the residues.  With two variables the whole
construction is carried out in $v$ with $w$ a symbol, over
$\Q(w,\eps)$ --- the singular points of that sliced equation are the
$v$-roots of the letters, algebraic functions of $w$; the freedom left
is a $v$-independent $U(w;\eps)$,
which must both keep the conjugated residues $U^{-1}R_aU$ constant
and satisfy the linear equation
$\partial_w U = A_w'U - \eps\,U\sum_a R_a''\,\partial_w\log\varphi_a$
--- a nullspace computation over $\Q(w,\eps)$ whose solvability is
decided by the solve itself and not by flatness, which is necessary
and not sufficient.  A candidate counts
as the canonical form only when the transformed block equals
$\eps\sum_a R_a\,\dd\log\varphi_a$ identically in \emph{both}
variables.  With
the diagonal canonical, the entries below it yield to unipotent
transformations $T = 1 + D$ (so $D^2 = 0$): writing $\eps e_k$ and
$\eps c_j$ for the canonical diagonal blocks of row $k$ and column
$j$, and $\bar b_{kj}$ for the coupling to be removed (including
the feed-through of previously completed entries of the same row),
the unknown block $D_{kj}$ and the residues $K_a$ of the completed
entry --- constant in the kinematics --- satisfy the
\emph{affine-linear} equation
\begin{equation}
  \dd D_{kj} = \eps( e_k D_{kj} - D_{kj} c_j ) + \bar b_{kj}
  - \eps \textstyle\sum_a K_a\, \dd\log\varphi_a .
\end{equation}
A linear operator applied to the unknowns plus the fixed source
$\bar b_{kj}$: its solution set is accordingly an affine subspace, one
particular solution plus the null space of the linear part, the same
structure the constraint system of the boundary analysis has below.
It is solved by an ansatz over the letters with
undetermined polynomial numerators, evaluated at many rational
kinematic and regulator values modulo large primes; the exact
rational solution --- a rational function of $\eps$ as well --- is
reassembled by the Chinese remainder theorem and rational
reconstruction, and proved by exact substitution.  Where the solution
is not unique, the representative with vanishing top-degree numerator
coefficients is chosen: all are correct, their canonical forms
differing by a constant transformation that preserves the diagonal
blocks, and a low-degree one keeps the degrees of the later rows from
compounding.  What survives is a dependence on $\eps$ that
is constant in the kinematics --- a residue of the $\eps$-dependent
normalization of the basis --- and it is removed last by a
conjugation $T(\eps)$, determined by the conditions
$\tilde R_a(\eps)\,T(\eps) = \eps\,T(\eps)\,R_a$, linear in $T$, at a
few rational kinematic points, and proved exactly on the full
connection.

\section{Square roots and their rationalization}
Eliminating the direction of an unobserved momentum at fixed
invariants requires solving quadratic equations, whose real solutions
end where a Gram determinant vanishes; for three momenta
$P_1, P_2, P_3 = -P_1-P_2$ that determinant is a K\"all\'en
polynomial,
\begin{equation}
  \det\begin{pmatrix}
    2P_1^2 & 2P_1\!\cdot\!P_2\\ 2P_1\!\cdot\!P_2 & 2P_2^2
  \end{pmatrix}
  = 4P_1^2P_2^2 - \big(P_3^2-P_1^2-P_2^2\big)^2
  = -\lambda\big(P_1^2,P_2^2,P_3^2\big) ,
\end{equation}
with $\lambda(a,b,c) = (c-a-b)^2-4ab$ the triangle function.
With the virtualities $s$, $-t$, $-u$ and $\lambda$ homogeneous of
degree two, this is, in units of $s^2$, the K\"all\'en polynomial
\begin{equation}
  \lambda_1 = \lambda(1,v,w) = (1-v-w)^2 - 4vw
  = \Big(1-\big(\sqrt v + \sqrt w\big)^2\Big)
    \Big(1-\big(\sqrt v - \sqrt w\big)^2\Big) ,
\end{equation}
whose second factor is positive on the whole triangle: the zero locus
is $\sqrt v + \sqrt w = 1$, a pseudo-threshold \emph{inside} the
physical region, through $v = w = \tfrac14$.  Alphabets of many
blocks therefore contain $\sqrt{\lambda_1}$-type letters.  Six
radicands occur in all; the involution $v \leftrightarrow w$
exchanges two pairs of them and fixes the other two.  The
\emph{root content} of a system is the number of independent square
roots in its alphabet, counted so that $\sqrt{q_1}$, $\sqrt{q_2}$,
$\sqrt{q_1 q_2}$ are two roots and not three --- formally, the rank
of the subgroup of $\Q(v,w)^\times/(\Q(v,w)^\times)^2$ that the
radicands generate, computed as the rank over $\mathbb{F}_2$ of the
matrix of exponent parities of the irreducible factors.  Reading the
radicands off the derived connection, from the discriminants of the
irreducible quadratic factors of its singular loci and of its residue
eigenvalues, gives an upper bound on that rank and not its value,
since an apparent singularity contributes a discriminant but no root.
In this channel about half the families have no root, a third have
one, a dozen have two, and three have three.  A single root is
removed by a \emph{rationalizing substitution}, a rational change of
variables whose Jacobian determinant is not identically zero and
under which the radicand becomes a perfect square: for instance
$v = xy$, $w = (1-x)(1-y)$ gives $\lambda_1 = (x-y)^2$, so
$\sqrt{\lambda_1} = x-y$ and the alphabet becomes rational in
$(x,y)$.  Such substitutions are the classical parametrization of a
conic by the pencil of lines through a rational point on it, here the
point at infinity: writing $q = \alpha^2 w^2 + b\,w + c$ as a
quadratic with square leading coefficient and setting
$\sqrt q = \alpha w + \tau$, the squares cancel and leave the linear
equation
\begin{equation}
  w = \frac{c - \tau^2}{2\alpha\tau - b} ,
  \qquad \sqrt q = \alpha\,w + \tau ,
\end{equation}
so that $(v,\tau)$ is a rationalizing chart.  Two roots are removed
by applying the same step to the second radicand \emph{in the
variables produced by the first}, which requires that it still be of
degree two there --- the pullback of a conic under a degree-two map
is in general a quartic --- and that its leading coefficient be a
square; for the K\"all\'en pairs both hold.

For three roots no rationalizing substitution exists.  Such a
substitution would be a dominant rational map from the plane onto the
$(\Z/2)^3$ cover of the plane branched over the three conics, a
complete intersection of three quadrics
$r_i^2 = Q_i(v,w,h) = h^2q_i(v/h,w/h)$ in $\mathbb{P}^5$, of degree
eight like the cover.  Its pairwise branch intersections are
transverse except for finitely many simple tangencies, each lifting to
two ordinary double points (nodes, locally $xy = z^2$); how many is
immaterial, since nodes are rational double points whose resolution is
crepant --- the canonical class pulls back without correction --- and
adjunction leaves it trivial, $\mathcal{O}(6-6)$, with vanishing
irregularity $h^1(\widetilde X, \mathcal{O}_{\widetilde X}) = 0$ and
geometric genus $p_g=1$: a K3 surface.  A unirational
surface, one dominated by a rational map from the plane, is rational
in characteristic zero by Castelnuovo's theorem, and a rational
surface has $p_g = 0$.  So one computes
instead in the extension $\Q(v,w,\eps)[r_1,r_2,r_3]/(r_i^2-q_i)$,
whose Galois group $(\Z/2)^3$ acts by sign changes.  Already for one
root the mechanism is complete: multiplication by $c_0 + c_1 r$ is
the matrix $\left(\begin{smallmatrix} c_0 & c_1q\\ c_1 & c_0
\end{smallmatrix}\right)$, differentiation obeys
$\dd r = \tfrac12 r\,\dd\log q$, and a one-form equation in the
algebra splits into two coupled equations with rational coefficients
for the even and odd components.  For three roots there are eight
graded components; the eight sign embeddings must and do recombine
consistently.

\section{Paths and polylogarithms}
With the canonical form established, evaluation means choosing a
base point in the physical triangle and a path composed of segments
on which every letter is linear in the path parameter.  The criterion
is what the generalized polylogarithms
\begin{equation}
  G(a_1,\dots,a_k;t) = \int_0^t \frac{\dd t'}{t'-a_1}\,
  G(a_2,\dots,a_k;t') , \qquad G(;t) = 1 ,
\end{equation}
integrated along a path from $0$ to $t$ avoiding the $a_i$ with a
fixed branch of each logarithm, and completed for the words in which
every $a_i$ vanishes --- the only ones the recursion cannot reach ---
by $G(\vec 0_k;t) = \log^k t/k!$, require: if
$\varphi_a(\gamma(t)) = \kappa_a\,(t-\alpha_a)$ then
$\dd\log\varphi_a = \dd t/(t-\alpha_a)$, the constant dropping out,
and every iterated integral along the segment is a $G$ with indices
the roots $\alpha_a$.  (Thus $G(1;t) = \log(1-t)$ and
$G(0,1;t) = -\operatorname{Li}_2(t)$, the two functions of the
example above.)  A letter quadratic on a segment stays quadratic on
every sub-segment, so one either chooses the segments differently ---
in practice parallel to the axes, where the bilinear letters of the
rationalizing substitutions become linear again --- or admits the
quadratic's linear factors as letters with algebraic constants.
Flatness makes the result depend only on the endpoint and the
homotopy class of the path; letters' zeros are crossed only as
deliberate analytic continuation.  The order in $\eps$ carried for
each master is fixed in advance: if the assembled cross section is
$\sigma = \sum_m c_m(\eps) I_m$, with $c_m$ carrying a pole of order
$p_m$, and $\sigma$ is wanted through $\eps^{\,n}$, then $I_m$ is
wanted through $\eps^{\,n+p_m}$.  Of the $347$ masters, $203$ are
needed only to their finite order and $131$ one order deeper, $9$
deeper still, and $4$ less deeply.

\section{Boundary conditions}
The general solution leaves one constant vector per family per
order, and almost all of it is fixed without evaluating any new
integral.  The most productive condition is the exclusion of local
behaviour.  Near $w = 0$ the equation offers
\begin{equation}
  \vec F(v,w) = \sum_\lambda w^{\lambda}\,
  \vec f_\lambda(v,\log w) ,
  \qquad
  \lambda \in \eps\operatorname{spec}R_0 + \Z_{\ge0} ,
\end{equation}
each $\vec f_\lambda$ linear in the boundary constants, while the
integral representation admits only exponents $n + a\eps$ produced by
its momentum regions.  The two lists must be compared in one basis:
the regions constrain the master $\vec I = T\vec F$, and its exponents
are those of the canonical $\vec F$ shifted by the order of $T$ at the
edge --- which is also where the negative integer parts used in the
next section come from.  The regulator enters a region only through the
measure, $\dd^dk \sim w^{4-2\eps}$ for a soft momentum $k \sim w$ and
$\dd^dk \sim w^{2-\eps}$ for a collinear one, so
$a = -2j_{\text{soft}} - j_{\text{coll}}$ counts the momenta a region
scales.  With two loop momenta,
$j_{\text{soft}} + j_{\text{coll}} \le 2$ and hence
$a \in \{0,-1,-2,-3,-4\}$: five values, which is the whole spectrum
observed on the soft edge.  Every offered but forbidden mode gives
$\vec f_\lambda \equiv 0$: one exact linear condition per mode ---
and the argument is the comparison of exponent lists, not an appeal
to finiteness, since $w^{-\eps}$ is itself bounded as $w\to0$ for
$\operatorname{Re}\eps<0$.  The rest follows from elementary
normalizations
(the phase-space volume is a closed ratio of $\Gamma$-functions and
fixes the lowest sector wherever it appears), from inheritance
(families are treated in subsector order, and every already-known
subsector master imposes its expansion wherever it recurs), and from
regularity at letters' zeros interior to the region of convergence.
All conditions assemble into one linear system $C\vec C^{(k)} =
\vec d^{\,(k)}$ per order, so the number of genuinely undetermined
constants is $\dim\ker C$, known before any evaluation.  After
identification across blocks --- the same subsector being the same
integral --- at most thirty-three candidate \emph{boundary periods}
remain for the whole channel, values of cut phase-space integrals at
ordered limits, $\zeta$-values and kin, established by exact
derivations with numerics as corroboration only.  Their character is
already visible in the beta integral
$\int_0^1\dd x\,x^{-1+a\eps}(1-x)^{-1+b\eps} =
\Gamma(a\eps)\Gamma(b\eps)/\Gamma((a+b)\eps)$, whose expansion is
$\frac{a+b}{ab\,\eps}\big[1 - ab\,\zeta_2\eps^2 + O(\eps^3)\big]$: a
$\zeta$-value out of $\Gamma$-functions alone.  Where the parametric
integrals do not factor, the Mellin--Barnes splitting
$(X+Y)^{-\lambda} = \frac{1}{\Gamma(\lambda)}\frac{1}{2\pi i}\int \dd
z\,\Gamma(\lambda+z)\Gamma(-z)X^z Y^{-\lambda-z}$ separates them
again.

\section{Edges of phase space}
Let $u$ be the distance to an edge of the triangle, rescaled to
$u\in[0,1]$: $u = 1-v-w$ on the soft edge, $u = v$ or $u = w$ on the
collinear ones.  The local solutions there are $u^{\,n+a\eps}(\log
u)^j$, with all data read from the residue matrix; a mode coefficient
is extracted from the global solution as
$\lim_{u\to0}\Pi_\lambda u^{-\eps R_u}\vec F$, the projection onto
the generalized eigenspace of exponent $\lambda$, and on the soft
edge it is a function of the variable surviving along that edge, not
a number.  The cross section then requires the distributional
identity
\begin{equation}
  u^{\,-1+a\eps} = \frac{1}{a\eps}\,\delta(u)
  + \sum_{k\ge0}\frac{(a\eps)^k}{k!}
  \Big[\frac{\log^k u}{u}\Big]_+ ,
  \qquad
  \int_0^1\!\dd u\,[g]_+f = \int_0^1\!\dd u\;g\,[f(u)-f(0)] .
\end{equation}
It is proved in one line: integrating against $f$ and adding and
subtracting $f(0)$ gives the term
$f(0)\int_0^1 u^{-1+a\eps}\,\dd u$, equal to $f(0)/(a\eps)$, plus a
remainder in
which $f(u)-f(0)$ vanishes at the endpoint, so that the expansion in
$\eps$ may be carried out term by term.  Modes $u^{\,n+a\eps}$ with
$n\le-2$ are treated by subtracting further Taylor terms and produce
derivatives of $\delta(u)$; for $n=-2$ the subtraction is
$f(u)-f(0)-u f'(0)$.  As written the identity is exact on the
collinear edges, where $u$ runs over $(0,1)$; on the soft edge the
rescaling is by $1-v$, so the $\delta$ coefficient carries
$(1-v)^{a\eps}$.  The identity exists only for the unexpanded
power: expanding the master in $\eps$ first leaves $u^{-1}\log^k u$,
non-integrable at the endpoint, and destroys the $\delta$-content.
The masters therefore carry their exact local mode data alongside the
polylogarithmic expansion, and the identity is applied before the
final expansion in $\eps$.

\section{Search and proof}
Nothing depends on how a transformation was found: an ansatz, a fit,
arithmetic modulo a prime produce candidates, and what makes a
candidate a theorem is an exact identity --- flatness for a
connection, a numerics-free derivation for a boundary period, and
for a canonical form
\begin{equation}
  T^{-1}A\,T - T^{-1}\dd T
  \;=\; \eps \sum_a R_a\,\dd\log\varphi_a ,
\end{equation}
verified as an identity of rational (or algebraic) matrices.  Where a
verification is done modulo a prime at random points, the bound is
stated: a nonzero polynomial of total degree $D$ vanishes at a
uniform point of $\Fp^{\,n}$ with probability at most $D/p$, so a
false identity survives $r$ independent evaluations with probability
at most $(D/p)^r$ --- below $10^{-29}$ already at $r = 2$ for
$D \le 10^4$ and a sixty-three--bit prime.  At the time of writing
$88$ of the $91$ families carry canonical forms established in this
way; the three three-root families await the Galois-algebraic
treatment, and the higher boundary periods and the local mode
analysis of every master at every edge remain open.  The method is
settled: each step is defined and each of its claims is decided by an
exact identity.  What remains is its execution on what is still open.

\end{document}
